%++++++++++++++++++++++++++++++++++++++++
% Don't modify this section unless you know what you're doing!
\documentclass[letterpaper,12pt]{article}
\usepackage{fullpage}
\renewcommand{\baselinestretch}{1.5} 
\usepackage{cite}
\usepackage{amsmath,amssymb,amsfonts}
\usepackage{algorithmic}
\usepackage{graphicx}
\usepackage{diagbox}
\usepackage[table]{xcolor}
\usepackage[margin=1in,letterpaper]{geometry} % decreases margins
\usepackage{textcomp}
\def\BibTeX{{\rm B\kern-.05em{\sc i\kern-.025em b}\kern-.08em
    T\kern-.1667em\lower.7ex\hbox{E}\kern-.125emX}}
\makeatletter
\let\NAT@parse\undefined
\makeatother
\usepackage{hyperref}
\hypersetup{
	pdftitle={PhD Thesis Proposal Hazoor Ahmad PhDEE17004},
   colorlinks=true,
   citecolor=blue,
   linkcolor=blue,
   urlcolor=blue
}
%++++++++++++++++++++++++++++++++++++++++
\begin{document}
\begin{titlepage}

\begin{center}
\Large \textbf {Hardware for Machine Learning}\\
\large \textbf {Subtitle if any}\\
\includegraphics[scale=0.7]{images/ITU_Logo.jpg}\\[0.1in]

\large \textbf{MS/PhD Thesis Proposal}\\[0.1in]

 \small
       \textit{\textbf{Submitted in partial fulfillment of
        the requirements for the award of the degree of}}\\[0.3in]
        
\normalsize
       \textbf{Master/DOCTOR OF PHILOSOPHY \\IN\\ ELECTRICAL ENGINEERING}\\[0.5in]

% Submitted by
By \\
\textbf{Hazoor Ahmad }\\ PhDEE17004




\vspace{.3in}
Supervisor Name \\
{\textbf{Dr. Rehan Hafiz}}\\[0.1in]
Co-Supervisor Name \\ 
{\textbf{Dr. Mohsin Ali}}\\[0.1in]

\vfill

% Bottom of the page
{Department of Electrical Engineering, Faculty of Engineering}\\
\normalsize
\textsc{Information Technology University}\\
Lahore, Punjab, Pakistan \\
August 2019

\end{center}

\end{titlepage}
\begin{abstract}

This section consists of the topic/research problem, theoretical approach, research methodology and significance of the study.
% Topic/Research Problem
ddasdad sggfgs fgfdsgsg.
% Theoretical Approach
ddasdad sggfgs fgfdsgsg.
% Research Methodology
ddasdad sggfgs fgfdsgsg.
% Significance
\end{abstract}
\pagebreak 
\section{Introduction} \label{intro}
This brief paragraph states the objective and/or goals of the PhD research. The information provided in this area gives the reader a quick overview of what is included in the proposal.\\
reseacha adaadad asdasdasd \cite{alexnet} fsdafh dsafas fasf safsa fasfas fas f.
\begin{figure}[htbp]
\centerline{\includegraphics[trim=0 0 0 100,clip,width=\columnwidth]{images/overview.pdf}}  % left bottom right top
\caption{Overview of Hardware for Machine Learning}
\label{overview}
\end{figure}
This research ais shown in Fig. ~\ref{overview}. Our focus is design strategies.\\
Organized as follows: section ~\ref{sota} contains a comprehensive literature survey on deep learning, section ~\ref{losota} limitations of state-of-the-art , section ~\ref{pspr} includes problem statement, its scope, novelty and justification, section ~\ref{significance} explains significance and impact of proposed approach, section ~\ref{preliminary-research} provides a detailed design, modelling, analysis, simulation and evaluation of preliminary research results, section ~\ref{plan} enlists a clear set tasks or major milestones and evaluation strategies with their time-line in the form of Gantt chart.

\pagebreak
\section{Literature Survey} \label{sota}
A detailed review of literature establishes your command in your area of research. This chapter should provide a complete and critical review of the state of the art work and not just summaries of books/articles.\\

This survey is focused on NeuFlow \cite{neuflow}. All these implementations include Zhang \cite{zhang}, Ayat \cite{ayat}, and Caffeine \cite{caffeine}.\\

% any division of survey classification 
gsdfgsdfgg gfsgsdf fdsgdgfsfdgsdf . 
 % Systolic Array Configurations
 SysArrAccel \cite{sysarraccel} pointed out exploration.
 SmartShuttle \cite{smartshuttle} provides configuration.
 Systimator \cite{systimator} have tested networks .\\
% optimization
Fused-Layer \cite{fused} performs t optize fo a ngle FPGA.\\
% Batch Processing and 3D systolic Array
Escher \cite{escher} utputs. 
.
\section{Limitations of the State of the Art} \label{losota}
In this section you need to clearly report the current Limitations of the State of the Art. The purpose is to clearly identify the limitations and weakness of the state of the art so that your chosen problem statement is justified.\\
State-of-the-art techniques lack \cite{sysarraccel,smartshuttle} or do not utilize \cite{escher,fused,maeri}. 
In addition to above, 3D systolic array has not yet been fully deployed for DNNs \cite{3d}. 
\pagebreak
\section{Problem Statement \& Proposed Research} \label{pspr}
In this section you need to report your identified Problem Statement. Furthermore, also provide a clear scope of your proposed research. The novelty of the problem being addressed and the particular approach being explored should also be commented upon/justified.\\
% Problem Statement
This research is focused on .
We will .
% Scope
The sessgdfg fgsdfgdf fdgsdg gsdg .
% Novelty
The sessgdfg fgsdfgdf fdgsdg gsdg .
% Justification
The sessgdfg fgsdfgdf fdgsdg gsdg .
\section{Significance of the Study} \label{significance}
This section is to provide the significance and impact of your proposed research. It should answer the following questions: Why you believe the study is significant? What implications your findings may have? Who will benefit from it? What will it contribute to the existing body of knowledge?\\
Systolic arrays are most useful .
\pagebreak 
\section{Preliminary Research and Results} \label{preliminary-research}
A written summary of some or all of the research performed is presented in a coherent manner. This section would include the approach taken and some preliminary results.\\


In this section we will demonstrate .
\begin{figure}[htbp]
\centerline{\includegraphics[trim=250 200 50 100,clip,width=\columnwidth]{images/figure1.pdf}} % left bottom right top
\caption{Dataflow engine.}
\label{syzerdataflow}
\end{figure}
\subsection{E-Syzer}
In this section we propose E-Syzer: An \underline{E}fficient \underline{Sy}stolic array si\underline{zer} for DNN . Our proposed Fig.~\ref{syzerdataflow}. 
We consider three different types of data traversals:
\begin{enumerate}
    \item  Input: 
    \item Weight: 
    \item Output: 
\end{enumerate}
\begin{figure}[htbp]
\centerline{\includegraphics[trim=50 200 50 100,clip,width=\columnwidth]{images/figure1.pdf}} % left bottom right top
\caption{Systolic array}
\label{fig2}
\end{figure}

systolic array in Fig.~\ref{syzerdataflow}.

\begin{figure*}\centering
\begin{tabular}{p{1.5in}p{1.5in}p{1.5in}p{1.5in}}
\includegraphics[width=1.5in]{images/tanh.eps} &
\includegraphics[width=1.5in]{images/dtanh.eps} &
\includegraphics[width=1.5in]{images/tanh.eps} &
\includegraphics[width=1.5in]{images/tanh.eps} \\
\tiny{(a) on($\rho=0$)} & 
\tiny{(b) ysis($\rho=0$)} &
\tiny{(c) Permance Ansis($\rho=0$)} &
\tiny{(d) Degn Spe($\rho=0$)}\\
\includegraphics[width=1.5in]{images/dtanh.eps} &
\includegraphics[width=1.5in]{images/tanh.eps} &
\includegraphics[width=1.5in]{images/dtanh.eps} &
\includegraphics[width=1.5in]{images/tanh.eps} \\
\tiny{(e) Layer wise resource estimation($\rho=1$)} & 
\tiny{(f) Resource Analysis($\rho=1$)} &
\tiny{(g) Performance Analysis($\rho=1$)} &
\tiny{(h) Design Space($\rho=1$)}
\end{tabular}
\caption{Layer rders.}
\label{fig3}
\end{figure*}

\subsubsection{Transr Enges (WTE, ITE, OTE)}
Weiht,  and oput tansfr gines.
\subsubsection{Systlic Aay (SA)}
Onevey clok cyle in Fig.~\ref{syzerdataflow}.

\begin{table}[htbp]
\caption{Syzer Parameters}
\begin{tabular}{|p{7cm}|p{3cm}|p{2.3cm}|p{2.3cm}|}
\hline \multicolumn{4}{|c|}{\textbf{CNN Parameters}} \\
\hline \textbf{Parameter Name} & \textbf{Symbol} & \textbf{Tiling} &  \textbf{indexing}\\
\hline output  &  $M$ & $T_m$ & $m$  \\
\hline  input  &  $N$ & $T_m$ & $n$ \\
\hline  rows &  $R$ & $T_m$ & $r$ \\
\hline columns &  $C$ & $T_m$ & $c$ \\
\hline \multicolumn{4}{|c|}{\textbf{FPGA/Hardware Resources}} \\
\hline No. of available DSP blocks  & $\mathbb{D}$ & &\\
\hline No. of available slice LUTs & $\mathbb{G}$ & &\\
\hline 
\end{tabular}
\label{tab1}
\end{table}

Before we move on, $\mathcal{O}$
\begin{equation}
T_m = min(\lceil\frac{T_m}{T_m}\rceil,R) \label{tr}
\end{equation}
For simplification we take all columns of a layer to be the tiling factor for ofm columns.
\begin{equation}
T_c = min(\lceil\frac{C_{TH}}{\mathcal{O}}\rceil,C)
\label{tc}
\end{equation}

To minimize optimization problem
\begin{equation*}
\begin{aligned}
& \underset{<\mathcal{A},\mathcal{B}>}{\text{minimize}}
& & \sum_{i=1}^{\mathcal{C}}{\mathfrak{D}_{i}} \\
& \text{subject to} \\
& & T_m \leq R\\
& & T_m \leq C\\
& & \mathfrak{d} \leq \mathbb{D}
\label{resourceopt}\end{aligned}
\end{equation*}

\subsubsection{Evaluation}


\pagebreak 
\section{Proposed Research Plan} \label{plan}
This objective of this section is to provide a clear set of tasks and intended approaches that shall be executed and evaluated to complete the PhD research work. All major milestones should be clearly mentioned. Evaluation plan can also be provided as to how the results shall be evaluated. A Gantt chart should be provided highlighting the outline of the work needed to complete the PhD research, and the time required for completion\\

There are following milestones to study
\begin{enumerate}
  \item First Goal
  \item Second Goal
  \item Third Goal
  \item Froth GOal
\end{enumerate}
MATLAB\textsuperscript{\tiny\textregistered} will be usd fr al simultions. Keras\textsuperscript{\tiny\textregistered}, TensorFlow\textsuperscript{\tiny\textregistered} or Caffe\textsuperscript{\tiny\textregistered}  fr taining. Harware implemntatn will be ing Vivado Desgn Sute - HLx Ediions by Xilinx\textsuperscript{\tiny\textregistered}.\\
Table ~\ref{gantt} provides a summary of major tasks along with their expected time of completion.
% Better make Gantt Chart as https://www.overleaf.com/latex/examples/an-example-of-gantt-chart-in-overleaf/xzbbbcfwdxhz
\begin{table}
\begin{tabular}{|p{4.5cm}||m{1.5cm}|m{1.5cm}|m{1.5cm}|m{1.5cm}|m{1.5cm}|m{1.5cm}|}
\hline \backslashbox[5cm]{Activity}{Time}&F 18&S 19&F 19&S 19&F 20&S 20\\
\hline 
\hline Literature Review and Study             &\cellcolor{gray} \begin{center}\Huge{\checkmark} \end{center} &&&&&\\
\hline Task 1                                  &&\cellcolor{gray} \begin{center}\Huge{\checkmark} \end{center} &&&&\\
\hline Task 2                                  &&&\cellcolor{gray} \begin{center}\Huge{\checkmark} \end{center} &&&\\
\hline Task 3                                  &&&&\cellcolor{gray} \begin{center}\Huge{\checkmark} \end{center} &&\\
\hline Task 4                                  &&&&&\cellcolor{gray} \begin{center}\Huge{\checkmark} \end{center} &\\
\hline Final Write-up \& Thesis Submission     &&&&&&\cellcolor{gray} \begin{center}\Huge{\checkmark} \end{center} \\
\hline
\end{tabular}
\caption{List of tasks}
\label{gantt}
\end{table}

\pagebreak 
\textbf{Other Considerations} \\
References should be presented in a consistent manner throughout the proposal. Use APA Referencing style for formatting the bibliographic material. It is important that figures, tables, and references in the proposal are presented in a manner consistent with professional publication standards. When placing a figure or table and its identifying description in the proposal, it is important to consider ease of access for the document reader. In most cases the text which introduces a figure or table will precede the placement of the figure or table in the proposal\\
\bibliography{Library_for_Automatic_References.bib} 
\bibliographystyle{apalike}
\end{document}